\clearpage
\section{Methodology}
\label{sec:methods}
\subsection{End-to-End Deep Learning Pipeline}
\subsubsection{Pipeline Overview}
The project separates configuration, data preparation, model construction, optimization, prediction, evaluation and checkpoint management~\cite{code}. The five models share this infrastructure. Training updates model weights, validation selects the checkpoint, and the final test only evaluates the selected model.
The processing stages are dataset preparation, saved splitting, training-only normalization, DataLoader construction, model forward passes, loss/optimization, validation-based checkpoint selection and held-out prediction/evaluation.

\subsubsection{DataLoader, Training and Evaluation}
Separate dataset objects provide independent training and validation transforms while referring to disjoint saved indices. The training loader shuffles with a seeded generator; validation and test retain their order. Batch size is 128 and \texttt{num\_workers=0}. Each training batch clears gradients, computes logits and cross-entropy, backpropagates and performs one optimizer step. Training enables dropout; validation and test disable it with evaluation mode and gradient suppression.

Epoch loss is sample-weighted: for batch size $n_b$ and mean batch loss $\ell_b$,
\begin{equation}
\bar\ell=\frac{\sum_b n_b\ell_b}{\sum_b n_b}.
\end{equation}
Thus the final partial batch is not weighted as heavily as a full batch. Accuracy and Macro-F1 are computed from the entire epoch's prediction vectors rather than averaged over batch scores. Evaluation rebuilds the architecture from the saved YAML, loads its selected checkpoint and exports test metrics. Error analysis separately verifies that its Accuracy and Macro-F1 agree with the saved test record.

\subsection{Architectures and Input Representations}
All models are trained from scratch. Table~\ref{tab:architectures} summarizes their representations; parameter counts include biases and are recorded in the branch's inference artifacts.
\begin{table}[!htbp]\centering\small
\caption{Architectures used in the main comparison. All outputs are ten logits.}
\label{tab:architectures}
\begin{tabularx}{\linewidth}{@{}lXr@{}}\toprule
Model & Representation and learned structure & Parameters\\\midrule
Linear & Flattened 784 pixels; one affine output layer & 7,850\\
MLP & $784\to256\to128\to10$; ReLU and dropout & 235,146\\
CNN & Two convolution/pooling blocks, then a 64-unit head & 220,234\\
LSTM & 28 row timesteps; two layers, hidden size 128 & 214,282\\
Transformer & 16 patch tokens; two 128-dimensional encoder layers & 274,954\\\bottomrule
\end{tabularx}\end{table}

\subsubsection{Linear / Softmax Classifier}
The baseline flattens the image to $v\in\mathbb R^{784}$ and computes
\begin{equation}
z=Wv+b,\qquad W\in\mathbb R^{10\times784},\quad b\in\mathbb R^{10}.
\end{equation}
It contains $784\times10+10=7{,}850$ parameters. Despite the name softmax classifier, the implementation returns logits directly. It learns a position-specific linear decision rule without hidden feature interactions, local weight sharing or explicit invariance. Its small parameter count makes it a useful cost baseline.

\subsubsection{Multilayer Perceptron}
The MLP maps $784\to256\to128\to10$. Each hidden layer has ReLU and dropout 0.2; the output layer has neither. Its total is
\begin{equation}
(784\cdot256+256)+(256\cdot128+128)+(128\cdot10+10)=235{,}146\;\text{parameters}.
\end{equation}
Nonlinear activations enable interactions between pixels, but flattening leaves spatial relations to fully connected weights. Dropout regularizes training and is disabled at evaluation.

\subsubsection{Convolutional Neural Network}
Two blocks apply $3\times3$ convolution with padding 1, ReLU and $2\times2$ max pooling. Channels grow from 1 to 32 to 64; feature maps shrink from $28\times28$ to $14\times14$ and then $7\times7$. The head flattens 3,136 features, maps them to 64 hidden units with ReLU and dropout 0.2, and outputs ten logits.
\begin{table}[!htbp]\centering\small
\caption{CNN learned layers. Pooling, ReLU, flattening and dropout add no learned parameters.}
\begin{tabularx}{\linewidth}{@{}lXr@{}}\toprule
Layer & Output / operation & Parameters\\\midrule
Convolution 1 & $32\times28\times28$; pool to $32\times14\times14$ & 320\\
Convolution 2 & $64\times14\times14$; pool to $64\times7\times7$ & 18,496\\
Hidden affine & $3{,}136\to64$ & 200,768\\
Output affine & $64\to10$ & 650\\\midrule
Total & & 220,234\\\bottomrule\end{tabularx}\end{table}
Local kernels and weight sharing encode the assumption that garment patterns are useful at multiple image locations. Pooling builds a coarser feature hierarchy. This structural bias differs from simply increasing dense-layer capacity.

\subsubsection{LSTM: Rows as Timesteps}
The channel dimension is removed to form $X\in\mathbb R^{B\times28\times28}$. There are 28 timesteps, each containing the 28 pixels of one image row. A two-layer, unidirectional LSTM with hidden size 128 processes the sequence from top to bottom. Its gates update a cell state and hidden state using the current row and preceding state. The final hidden state of the last layer summarizes the image; dropout 0.2 and a $128\to10$ affine layer produce logits. Inter-layer LSTM dropout is also 0.2.

For one layer, PyTorch stores four input-gate weight blocks, four recurrent blocks and two bias vectors per gate. Thus the count is
\begin{equation}
\underbrace{4(128)(28+128+2)}_{80{,}896}
+\underbrace{4(128)(128+128+2)}_{132{,}096}
+\underbrace{128\cdot10+10}_{1{,}290}=214{,}282\;\text{parameters}.
\end{equation}
The model can combine information across distant rows, but its sequential computation and row ordering do not explicitly encode a two-dimensional local neighborhood.

\subsubsection{Transformer: Patch Tokens and Self-Attention}
A $7\times7$ convolution with stride 7 projects each non-overlapping image patch to 128 features. The $4\times4$ patch grid becomes 16 tokens, $X\in\mathbb R^{B\times16\times128}$. A learned $1\times16\times128$ positional embedding distinguishes token locations. Each of two pre-normalized encoder layers uses four attention heads of width 32, residual connections, dropout 0.2 and a GELU feed-forward network of width 256. For each head,
\begin{equation}
Q=XW_Q,\quad K=XW_K,\quad V=XW_V,\qquad
\operatorname{Attention}(Q,K,V)=\operatorname{softmax}\!\left(\frac{QK^\top}{\sqrt{32}}\right)V.
\end{equation}
Queries, keys and values all come from the image-token sequence. Attention mixes information between patch locations; the combined encoder output retains shape $B\times16\times128$. Mean pooling over tokens, LayerNorm and a $128\to10$ head yield logits. There is no class token or pretrained backbone.

Patch projection contributes 6,400 parameters, positional embeddings 2,048, the two encoder layers 264,960, and the final normalization/head 1,546: 274,954 in total. Attention provides global token interactions but has less built-in local image structure than the CNN. The $7\times7$ patch size also changes how fine spatial details enter the network.

\subsection{Loss and Evaluation Protocol}
\subsubsection{Cross-Entropy and Prediction}
All models minimize unweighted cross-entropy on raw logits~\cite{pytorchloss}:
\begin{equation}
\mathcal L=-\frac1B\sum_{i=1}^{B}\log\frac{\exp z_{i,y_i}}{\sum_{c=0}^{9}\exp z_{ic}}.
\end{equation}
The predicted class is $\hat y_i=\arg\max_c z_{ic}$. Applying softmax before this loss is unnecessary; no additional threshold or post-processing is used.

\subsubsection{Accuracy, Macro-F1 and Class Diagnostics}
Accuracy is $N^{-1}\sum_i\mathbf1(\hat y_i=y_i)$. For class $c$,
\begin{equation}
F1_c=\frac{2TP_c}{2TP_c+FP_c+FN_c},\qquad
\mathrm{Macro\mbox{-}F1}=\frac1{10}\sum_{c=0}^{9}F1_c.
\end{equation}
The metric code explicitly includes all ten labels and uses zero for undefined denominators. Confusion-matrix rows are true classes and columns are predictions. Row-normalized diagonals measure recall, not overall accuracy. Saved error-analysis records include each class's precision, recall, F1 and support, plus the most frequent directed confusions.

\subsection{Training Setup and Experimental Design}
\subsubsection{Hardware, Dependencies and Randomness}
The branch records Windows x64, an Intel Core i7-8665U (four cores/eight threads), 16 GB RAM and CPU execution. Its declared environment is Python 3.12.10, PyTorch 2.14.0+cpu, torchvision 0.29.0+cpu, NumPy 2.5.3, scikit-learn 1.9.1, Matplotlib 3.11.2, seaborn 0.13.2, pytest 9.1.1, PyYAML 6.0.3 and Pillow 12.3.0. This is the source branch's environment record, not the current report-authoring machine's hardware.

All saved configurations use seed 42 and deterministic mode. The implementation seeds Python, NumPy and PyTorch, configures deterministic algorithms and seeds loader generators/workers. These controls do not guarantee bitwise equality across platforms or library versions~\cite{pytorchrepro}. No independent seed repeats are present in the main comparison.

\subsubsection{Shared Hyperparameters and Regularization}
\begin{table}[!htbp]\centering
\caption{Resolved settings shared by all five reported runs.}
\label{tab:training}
\begin{tabularx}{\linewidth}{@{}lX@{}}\toprule
Setting & Value\\\midrule
Optimizer / learning rate & Adam / $5\times10^{-4}$\\
Batch size / epochs & 128 / 20\\
Weight decay & 0\\
Seed / workers & 42 / 0\\
Train / validation / test & 51,000 / 9,000 / 10,000\\
Augmentation & Disabled\\
Checkpoint selection & Highest validation Macro-F1\\
Scheduler / early stopping & Neither used\\
Precision & Default float32; no mixed precision\\\bottomrule
\end{tabularx}\end{table}
Dropout is 0.2 in the nonlinear models and absent from Linear. The reported runs use one baseline configuration for each architecture. Equal epochs do not imply equal compute or matched capacity.

\subsubsection{Checkpoint Selection and Comparison Criteria}
After each epoch, a strict improvement in validation Macro-F1 replaces the checkpoint; an exact tie retains the earlier epoch. Checkpoints save model and optimizer states, epoch, validation metrics and configuration. Test evaluation uses that checkpoint rather than necessarily epoch 20.

The changed factors are architecture and input representation (flattened pixels, spatial feature maps, row sequence or patch tokens). The fixed factors are dataset split, normalization, seed, optimizer, learning rate, batch size, epoch budget, selection criterion and test metrics. The main hypotheses are that nonlinear interactions improve over a linear rule, and that convolutional locality offers a useful accuracy--cost trade-off. Evidence is evaluated through test Accuracy and Macro-F1 together with errors, parameters and elapsed time. This design compares concrete configurations; it cannot isolate representation from architecture or prove a universal ranking.

\subsubsection{Timing Protocol}
Training time is wall-clock duration of the fit loop, including training, validation, metric computation and checkpoint writes; dataset preparation and model construction are excluded. Inference artifacts record CPU execution, four threads, batch size one, five warm-up passes and 20 timed repeats using \texttt{perf\_counter}. Only the loaded model's forward pass is timed in inference mode. Data loading, preprocessing, checkpoint loading and startup are excluded. The report gives the median; p95 is retained for context. These are repeated timings of one preprocessed test image, not a distribution over the entire test set or end-to-end deployment latency.

\subsection{Implementation and Verification}
The group defines architecture assembly, configuration handling, normalization, split validation, training/evaluation loops, checkpointing, error analysis and reporting. PyTorch supplies layers, autograd, optimizers and loaders; torchvision supplies Fashion-MNIST/transforms; scikit-learn supplies classification metrics; Matplotlib/seaborn supply plots and PyYAML reads configurations. No pretrained model replaces any of the five mandatory architectures.

The updated branch collects pytest tests for metrics, weighted losses, training updates, evaluation behavior, checkpoint restoration, deterministic splitting, loaders, transforms, EDA and random-number control. Dataset tests replace downloads with in-memory inputs. These tests cover implementation behavior. For this report, the saved comparison CSV was independently checked against every selected run's JSON records, all checkpoint-epoch fields were reconciled, and Accuracy/Macro-F1 were recomputed from all five stored confusion matrices.

\subsubsection{Reproducibility}
\begin{itemize}
\setlength{\itemsep}{3pt}\setlength{\parskip}{0pt}
\item \textbf{Source and commands.} The experiment snapshot is \texttt{a1} commit \texttt{27ef182}~\cite{code}. Its \href{https://github.com/Sunnyngooo/CO3133_Group05_Assignment/blob/27ef182308f0f6ac6f04510ee9fc77b9a16cc14a/assignment1/README.md}{Assignment 1 README} provides environment setup, dataset preparation, training, evaluation and diagnostic commands.
\item \textbf{Saved records.} For each run listed in Section~\ref{sec:results}, \texttt{assignment1/runs/<run-id>/} stores \texttt{config.yaml}, \texttt{history.json}, \texttt{test\_metrics.json}, \texttt{inference\_metrics.json} and \texttt{error\_analysis.json}. These records preserve preprocessing settings, selected epoch and reported metrics.
\item \textbf{Checkpoint reconstruction.} Load \texttt{assignment1/checkpoints/<run-id>/best.pt} with the corresponding saved configuration. The checkpoint, history and metric files agree on the validation-selected epoch. Portable repository paths replace the original machine's absolute paths recorded in histories.
\end{itemize}
